\NSPage{151}%
\begin{EESegment}{EE-TGT-P151-001}
uniformly in \NSi{NS-F-P151-001}{\mathsf C,t_1}. Its real value stays bounded above and below. It
follows that \NSi{NS-F-P151-002}{\mathsf C E_{\mathrm r}=\sqrt{2X}\phi_{\mathrm r}} and its
reciprocal have the stated bounds away from \NSi{NS-F-P151-003}{X=0}. Also,
\end{EESegment}
\NSFormula{NS-F-P151-004}{display}{B.24}
\begin{equation}
  \Pi_{\mathrm r}(X)-\Pi_0=\mathsf C^{-2}\int_0^X\phi_{\mathrm r}(X',\eta)^2\,dX',
  \qquad
  \sup_{X\leq X_i}\lVert\Pi_{\mathrm r}-\Pi_0\rVert_{C_\eta^k}\leq C_k\mathsf C^{-2}.
\tag{B.24}\label{eq:B.24}
\end{equation}
\begin{EESegment}{EE-TGT-P151-002}
The regular axis factor makes this integral finite at zero.
\end{EESegment}

\begin{EESegment}{EE-TGT-P151-003}
For the sign of the source, the leading gradient term is
\NSi{NS-F-P151-005}{-H_*\widetilde\xi_0=L\Lambda\chi}. The bound
\NSi{NS-F-P151-006}{|\widetilde\xi_0|\leq L\Lambda\sigma_*^{-1}\sqrt\chi} shows that replacing
\NSi{NS-F-P151-007}{H_*} by \NSi{NS-F-P151-008}{H_{c,\mathrm r}} creates an error bounded by
\NSi{NS-F-P151-009}{C_{\sigma_*}\sqrt\chi+o_{t_1}(1)} after
\NSi{NS-F-P151-010}{\Lambda,\mathsf C} have been fixed. The logarithmic slope of the axis profile,
together with the remaining gradient, contributes \NSi{NS-F-P151-011}{C\chi+C/\Lambda}. Absorb
these errors into a small fraction of \NSi{NS-F-P151-012}{L\Lambda\chi}, using
\NSi{NS-F-P151-013}{C_{\sigma_*}\sqrt\chi\leq\varepsilon\Lambda L\chi+C_{\varepsilon,\sigma_*}/\Lambda}.
The constant term keeps the room provided by \NSi{NS-F-P151-014}{-W_*>2.8}, which gives the first
inequality in \eqref{eq:B.23}. The axis profile satisfies \NSi{NS-F-P151-015}{p_1=a>0}, so the
logarithmic \NSi{NS-F-P151-016}{\phi} slope is nonpositive. Cutting this slope off at zero gives
\NSi{NS-F-P151-017}{l_{\mathrm r}\leq1}.
\end{EESegment}

\begin{EESegment}{EE-TGT-P151-004}
The axial source satisfies
\NSi{NS-F-P151-018}{S_{n,\mathrm r}=Z_*+O(\Lambda^{-1})+o_{t_1}(1)+O(\mathsf C^{-2})}. On the
short cutoff, \NSi{NS-F-P151-019}{\mathcal D_XU_{\mathrm{nat}}=O(\Lambda^{-1})}. On the whole
finite interval, \eqref{eq:B.24} controls the pressure and the parameter derivative of the pressure
that is needed here. The equations for \NSi{NS-F-P151-020}{p_{1,\mathrm r},n_{s,\mathrm r}} are
\end{EESegment}
\NSFormula{NS-F-P151-021}{display}{B.25}
\begin{equation}
  \mathcal D_Xp_{1,\mathrm r}=\frac{XS_{q,\mathrm r}}{L}-l_{\mathrm r}p_{1,\mathrm r},
  \qquad
  \mathcal D_Xn_{s,\mathrm r}+n_{s,\mathrm r}=\frac{S_{n,\mathrm r}}{L}.
\tag{B.25}\label{eq:B.25}
\end{equation}
\begin{EESegment}{EE-TGT-P151-005}
The regular initial values and bounded coefficients give the stated parameter bounds. In the region
\NSi{NS-F-P151-022}{\chi\leq.99}, the axis alternative
\NSi{NS-F-P151-023}{|Z_*|>\delta_*} keeps the sign fixed and keeps
\NSi{NS-F-P151-024}{|n_{s,\mathrm r}|} above a positive lower bound.
\end{EESegment}

\begin{EESegment}{EE-TGT-P151-006}
The first equation in \eqref{eq:B.25} preserves positivity and gives both comparisons below.
\end{EESegment}
\NSFormula{NS-F-P151-025}{display}{none}
\[
\begin{aligned}
  p_{1,\mathrm r}(y)&\geq p_{1,\mathrm r}(0)e^{-y}+1.88\chi(e^y-e^{-y}),\\
  p_{1,\mathrm r}(X)&\geq p_{1,\mathrm r}(X_0)\frac{X_0}{X}+1.2\left(X-\frac{X_0^2}{X}\right).
\end{aligned}
\]
\begin{EESegment}{EE-TGT-P151-007}
In the region \NSi{NS-F-P151-026}{\chi>.99}, the first right-hand side stays above
\NSi{NS-F-P151-027}{2.3}. This is because its initial value can be replaced by
\NSi{NS-F-P151-028}{2.3}, and \NSi{NS-F-P151-029}{1.88\chi>2.3/2}. In the other region,
\NSi{NS-F-P151-030}{p_{2,\mathrm r}=Xn_{s,\mathrm r}/E_{\mathrm r}} grows uniformly in absolute
value with \NSi{NS-F-P151-031}{\mathsf C}, while \NSi{NS-F-P151-032}{p_{1,\mathrm r}} stays
positive and bounded. This proves \NSi{NS-F-P151-033}{v_{\mathrm r}>2+c}. The second comparison
gives \NSi{NS-F-P151-034}{p_{1,\mathrm r}>3} from \NSi{NS-F-P151-035}{X=100} onward.
\end{EESegment}
\end{proof}

\begin{EESegment}{EE-TGT-P151-008}
\subsection{Continuing to a nonzero leading residual stress.}\label{sec:B.6}
This subsection uses the reference bounds from Lemma~\ref{lem:B.4} to reduce the shear
\NSi{NS-F-P151-036}{\mathfrak s}, while keeping the integrated residual coefficients
\NSi{NS-F-P151-037}{p_s} close to their reference values as stated in \eqref{eq:B.28}. The stress
\NSi{NS-F-P151-038}{\mathcal T_0=F(p_s-\mathfrak s)} from \eqref{eq:4.11} must turn on smoothly
at the end of the axis region. It must satisfy the admissible stress cone on an initial collar next
to that region, and it must keep satisfying the relaxed cone condition throughout the rest of the
continuation. Proposition~\ref{prop:B.5} does this.
\end{EESegment}

\begin{proposition}\label{prop:B.5}
\begin{EESegment}{EE-TGT-P151-009}
The analytic axis profile can be continued to \NSi{NS-F-P151-039}{X_i} while keeping
\NSi{NS-F-P151-040}{E} positive and leaving the profile unchanged for
\NSi{NS-F-P151-041}{X\leq X_0}. Its stress factors as
\NSi{NS-F-P151-042}{\mathcal T_0=e_aB_0}. The scalar factor \NSi{NS-F-P151-043}{e_a} vanishes to
infinite order at \NSi{NS-F-P151-044}{X_0}, while the smooth vector coefficient
\NSi{NS-F-P151-045}{B_0} is nonzero there. The admissible stress cone holds on an initial collar,
and the strict relaxed cone condition holds on the rest of
\NSi{NS-F-P151-046}{(X_0,X_i]}. At \NSi{NS-F-P151-047}{X_i},
\NSi{NS-F-P151-048}{a=.8}, \NSi{NS-F-P151-049}{\mathcal D_XU=0}, and
\NSi{NS-F-P151-050}{p_1>2}.
\end{EESegment}
\end{proposition}
\begin{proof}
\begin{EESegment}{EE-TGT-P151-010}
First reduce the reference shear on a short interval. During that reduction,
\NSi{NS-F-P151-051}{p_s} changes by a small error that is kept under control. The residual stress
then becomes nonzero and satisfies the cone conditions just stated. Next, show that the stress has a
factor that vanishes to infinite order, and finish the continuation with small shear.
\end{EESegment}

\NSPage{152}%
\begin{EESegment}{EE-TGT-P152-001}
Fix a sufficiently large finite \NSi{NS-F-P152-001}{\mathsf C}. Choose
\NSi{NS-F-P152-002}{0<t_*\leq\bar t} small enough that the bounds in Lemma~\ref{lem:B.4} hold
uniformly for every reference width \NSi{NS-F-P152-003}{0<t_1\leq t_*}. The parameter bounds in
that lemma, along with its positive lower comparison for \NSi{NS-F-P152-004}{p_{1,\mathrm r}}, give
one finite upper bound \NSi{NS-F-P152-005}{V_{\max}} for
\NSi{NS-F-P152-006}{v_{\mathrm r}} across this whole family of reference profiles. This bound may
depend on \NSi{NS-F-P152-007}{\mathsf C}, but it does not depend on
\NSi{NS-F-P152-008}{t_1}. Choose \NSi{NS-F-P152-009}{0<\kappa_0<1/2} later from this bound. On
\NSi{NS-F-P152-010}{0<y<t_1}, set
\end{EESegment}
\NSFormula{NS-F-P152-011}{display}{B.26}
\begin{equation}
  e_a=(1-\kappa_0)\sigma(y/t_1),\quad
  \kappa=1-e_a,\quad
  a=\kappa p_{1,\mathrm r},\quad
  \mathcal D_XU=-\frac{\kappa Xn_{s,\mathrm r}}{2},\quad
  \partial_y\log\phi=-\frac{\kappa p_{1,\mathrm r}}{2}.
\tag{B.26}\label{eq:B.26}
\end{equation}
\begin{EESegment}{EE-TGT-P152-002}
On this interval, the reference profile agrees with the analytic axis profile. Subtracting the
reference equations gives the exact identities
\end{EESegment}
\NSFormula{NS-F-P152-012}{display}{B.27}
\begin{equation}
  \partial_y(\log\phi-\log\phi_{\mathrm r})=\frac{e_ap_{1,\mathrm r}}{2},
  \qquad
  \partial_y(U-U_{\mathrm r})=\frac{e_aXn_{s,\mathrm r}}{2}.
\tag{B.27}\label{eq:B.27}
\end{equation}
\begin{EESegment}{EE-TGT-P152-003}
Since \NSi{NS-F-P152-013}{e_a} is increasing, every fixed parameter derivative of these field
differences is \NSi{NS-F-P152-014}{O(ye_a)}. Constants in the activation comparisons may depend on
the already fixed \NSi{NS-F-P152-015}{\mathsf C,\Lambda}, but they stay uniform as
\NSi{NS-F-P152-016}{t_1} and \NSi{NS-F-P152-017}{\kappa_0} decrease within their allowed ranges.
The bounds for the reference family control the integrated coefficients independently of both
widths. Apply the radial moment identities \eqref{eq:4.16}, using one more parameter derivative.
They give
\end{EESegment}
\NSFormula{NS-F-P152-018}{display}{B.28}
\begin{equation}
  p_s-p_{s,\mathrm r}=O(ye_a),
  \qquad
  E_{\mathrm r}/E=1+O(ye_a).
\tag{B.28}\label{eq:B.28}
\end{equation}
\begin{EESegment}{EE-TGT-P152-004}
This comparison holds for the field values and their parameter derivatives. Radial derivatives of
the narrow cutoffs may be large.
\end{EESegment}

\begin{EESegment}{EE-TGT-P152-005}
The actual shear is
\NSi{NS-F-P152-019}{\kappa(p_{1,\mathrm r},p_{2,\mathrm r}E_{\mathrm r}/E)}. The common factor
\NSi{NS-F-P152-020}{\kappa} cancels from the ratio \NSi{NS-F-P152-021}{t_s=-b_s/a} before that
ratio is estimated. The definitions of \NSi{NS-F-P152-022}{P_{\mathrm c},J_{\mathrm c},v_s}
give
\end{EESegment}
\NSFormula{NS-F-P152-023}{display}{B.29}
\begin{equation}
  t_s=\frac{p_{2,\mathrm r}E_{\mathrm r}}{p_{1,\mathrm r}E},\qquad
  v_s=\kappa v_{\mathrm r}+O(ye_a),\qquad
  P_{\mathrm c}=v_{\mathrm r}+O(ye_a),\qquad
  J_{\mathrm c}=O(ye_a).
\tag{B.29}\label{eq:B.29}
\end{equation}
\begin{EESegment}{EE-TGT-P152-006}
No inverse power of \NSi{NS-F-P152-024}{\kappa_0} appears in these error constants. If
\NSi{NS-F-P152-025}{t_1} is small enough, then
\NSi{NS-F-P152-026}{P_{\mathrm c}-v_s\geq ce_a},
\NSi{NS-F-P152-027}{P_{\mathrm c}>2+c}, and
\end{EESegment}
\NSFormula{NS-F-P152-028}{display}{none}
\[
  (v_s-2)_+J_{\mathrm c}^2<2(P_{\mathrm c}-v_s)^2
  \qquad(0<y\leq t_1).
\]
\begin{EESegment}{EE-TGT-P152-007}
The ratio of the left side of this inequality to
\NSi{NS-F-P152-029}{(P_{\mathrm c}-v_s)^2} is at most \NSi{NS-F-P152-030}{Cy^2}. When
\NSi{NS-F-P152-031}{v_s>2}, the quadratic test in Lemma~\ref{lem:4.5} proves the admissible stress
cone condition. When \NSi{NS-F-P152-032}{v_s\leq2}, the inequality
\NSi{NS-F-P152-033}{P_{\mathrm c}>2} directly gives the strict relaxed condition. Since
\NSi{NS-F-P152-034}{v_{\mathrm r}(0)>2+c}, continuity gives a first collar of positive width where
\NSi{NS-F-P152-035}{v_s>2+c/2}. For the final fixed parameters, the collar is uniform in
\NSi{NS-F-P152-036}{\eta}.
\end{EESegment}

\begin{EESegment}{EE-TGT-P152-008}
The factor \NSi{NS-F-P152-037}{e_a=e^{-t_1^2/y^2}g_a(y)} is flat at zero, meaning that it vanishes
there to every order, and \NSi{NS-F-P152-038}{g_a} is smooth and positive. To show that division
by this flat factor still gives a smooth function, apply Lemma~\ref{lem:A.9} with
\NSi{NS-F-P152-039}{c=t_1^2}, \NSi{NS-F-P152-040}{j=0}, and coefficient
\NSi{NS-F-P152-041}{g_a b}. For every smooth coefficient \NSi{NS-F-P152-042}{b}, the lemma gives
\end{EESegment}
\NSFormula{NS-F-P152-043}{display}{none}
\[
  \frac{1}{e_a(y)}\int_0^y e_a(u)b(u,\eta)\,du
  =y^3B(y,\eta),
  \qquad B\in C^\infty.
\]
\begin{EESegment}{EE-TGT-P152-009}
Dividing by the positive smooth function \NSi{NS-F-P152-044}{g_a(y)} keeps the result smooth and
keeps every fixed order of mixed derivatives bounded. The differences in \eqref{eq:B.27} lie in
\NSi{NS-F-P152-045}{e_ay^3C^\infty}. Products, smooth compositions, and reciprocals of fields that
never vanish stay in this class. The moment and pressure differences vanish before
\NSi{NS-F-P152-046}{X_0}. One more integration puts them in
\NSi{NS-F-P152-047}{e_ay^6C^\infty}. Substituting into \eqref{eq:4.16} then proves
\eqref{eq:B.28} after division by \NSi{NS-F-P152-048}{e_a}, for every fixed mixed derivative. It
also gives the exact factorization
\end{EESegment}
\NSFormula{NS-F-P152-049}{display}{B.30}
\begin{equation}
  \mathcal T_0=e_aB_0(y,\eta),
  \qquad
  B_0(0,\eta)=F(X_0,\eta)p_{s,\mathrm r}(X_0,\eta)\neq0,
  \qquad
  B_0\in C^\infty.
\tag{B.30}\label{eq:B.30}
\end{equation}
\NSPage{153}%
\begin{EESegment}{EE-TGT-P153-001}
At the edge, the limiting shear is \NSi{NS-F-P153-001}{p_{s,\mathrm r}}. This gives
\NSi{NS-F-P153-002}{B_0\cdot(-t_s,1)=0} and
\NSi{NS-F-P153-003}{B_0\cdot(1,t_s)>0} there. At the endpoint,
\NSi{NS-F-P153-004}{B_0\cdot(1,t_s)>0} and
\NSi{NS-F-P153-005}{2(B_0\cdot(1,t_s))^2-(v_s-2)(B_0\cdot(-t_s,1))^2>0}. Both quantities stay
above positive lower bounds on a smaller collar. In particular, for every fixed radial and parameter
multi-index \NSi{NS-F-P153-006}{I},
\end{EESegment}
\NSFormula{NS-F-P153-007}{display}{B.31}
\begin{equation}
  |\mathcal T_0|\geq ce^{-t_1^2/y^2},
  \qquad
  |\partial^I\mathcal T_0|\leq C_Ie^{-t_1^2/y^2}y^{-N_I}.
\tag{B.31}\label{eq:B.31}
\end{equation}
\begin{EESegment}{EE-TGT-P153-002}
The derivatives in \eqref{eq:B.31} are taken in \NSi{NS-F-P153-008}{(y,\eta)}. Any fixed number
of derivatives in \NSi{NS-F-P153-009}{X} has the same form because
\NSi{NS-F-P153-010}{X_0>0}. The nonzero vector
\NSi{NS-F-P153-011}{B_0(0,\eta)} gives the limiting direction of the stress.
\end{EESegment}

\begin{EESegment}{EE-TGT-P153-003}
Continue the profile to \NSi{NS-F-P153-012}{X_i} while keeping the relaxed cone condition. After
the first transition, keep the prescription \eqref{eq:B.26} with
\NSi{NS-F-P153-013}{\kappa=\kappa_0} through the reference cutoff and the interval of constant
slope. Once \NSi{NS-F-P153-014}{\mathsf C,\Lambda} are fixed, integration over
\NSi{NS-F-P153-015}{[X_0,X_{\mathrm b}]} and \eqref{eq:4.16} show that the field values, their
logarithms, and the required derivatives differ by \NSi{NS-F-P153-016}{O(t_1+\kappa_0)} when
those derivatives are taken with respect to \NSi{NS-F-P153-017}{\eta}. Choose
\NSi{NS-F-P153-018}{\kappa_0} small compared with
\NSi{NS-F-P153-019}{V_{\max}} and the uniform comparison constants for the reference family. Then
choose \NSi{NS-F-P153-020}{t_1} small. This order does not require fixing a reference width before
\NSi{NS-F-P153-021}{\kappa_0}. Formula \eqref{eq:B.29} now gives
\NSi{NS-F-P153-022}{P_{\mathrm c}>2+c/2} and \NSi{NS-F-P153-023}{v_s<1}. The strict relaxed cone
condition follows no matter how large \NSi{NS-F-P153-024}{J_{\mathrm c}} is.
\end{EESegment}

\begin{EESegment}{EE-TGT-P153-004}
At \NSi{NS-F-P153-025}{X_{\mathrm b}}, multiply the prescribed value of
\NSi{NS-F-P153-026}{\mathcal D_XU} by a smooth factor \NSi{NS-F-P153-027}{\beta} that decreases
from one to zero over a short logarithmic interval. During this transition,
\NSi{NS-F-P153-028}{P_{\mathrm c}=p_{1,\mathrm r}+\beta p_{2,\mathrm r}^2/p_{1,\mathrm r}+o(1)>2},
while \NSi{NS-F-P153-029}{v_s<1}. Next, use a convex interpolation to move
\NSi{NS-F-P153-030}{a} from \NSi{NS-F-P153-031}{\kappa_0p_{1,\mathrm r}} to
\NSi{NS-F-P153-032}{.8}, keeping \NSi{NS-F-P153-033}{\mathcal D_XU=0}. First arrange that
\NSi{NS-F-P153-034}{\kappa_0\sup p_{1,\mathrm r}<.8}. During this second transition,
\NSi{NS-F-P153-035}{v_s=a\leq.8}. Also, \NSi{NS-F-P153-036}{P_{\mathrm c}=p_1>2}, by continuity
from the positive margin at \NSi{NS-F-P153-037}{X_{\mathrm b}}. Both transitions end strictly
before \NSi{NS-F-P153-038}{X_i}.
\end{EESegment}

\begin{EESegment}{EE-TGT-P153-005}
On the last interval of constant slope, \NSi{NS-F-P153-039}{a=.8} and
\NSi{NS-F-P153-040}{l=.6}. The same absorption of the gradient used in Lemma~\ref{lem:B.4} gives
\NSi{NS-F-P153-041}{S_q>1}. Its constant term is
\NSi{NS-F-P153-042}{.6(-W_*)>1.68}, while its leading gradient is
\NSi{NS-F-P153-043}{L\Lambda\chi}. Suppose that \NSi{NS-F-P153-044}{p_1=2} crossed downward.
At such a crossing,
\NSi{NS-F-P153-045}{\mathcal D_Xp_1=XS_q/L-.6p_1>X-1.2>0}, which is impossible. This finishes
the continuation.
\end{EESegment}
\end{proof}

\begin{EESegment}{EE-TGT-P153-006}
The continuation is still analytic in \NSi{NS-F-P153-046}{\eta} on the first collar. Reserve a
smaller rectangle inside that collar. Later changes will then leave both this regularity and the
factorization of the inner stress untouched.
\end{EESegment}

\begin{corollary}\label{cor:B.6}
\begin{EESegment}{EE-TGT-P153-007}
There is a \NSi{NS-F-P153-047}{t_c>0} such that \NSi{NS-F-P153-048}{t_c<t_1} and
\NSi{NS-F-P153-049}{4e^{t_c}<4.1}. The functions
\NSi{NS-F-P153-050}{E/\sqrt{2X},U,V_0/X,\Pi} are defined on
\NSi{NS-F-P153-051}{[0,X_0e^{t_c}]\times[-1,1]}, and they are smooth in
\NSi{NS-F-P153-052}{(X,\eta)} through \NSi{NS-F-P153-053}{X=0}, and they are analytic in
\NSi{NS-F-P153-054}{\eta} on one complex neighborhood. Every fixed radial derivative of these four
functions is analytic on that same parameter neighborhood and is bounded on a smaller one. On
\NSi{NS-F-P153-055}{X_0<X\leq X_0e^{t_c}}, the admissible stress cone condition and the
factorization \eqref{eq:B.30} hold. Every later change to the profile is supported strictly to the
right of \NSi{NS-F-P153-056}{X_0e^{t_c}}.
\end{EESegment}
\end{corollary}
\begin{proof}
\begin{EESegment}{EE-TGT-P153-008}
Choose a compact complex neighborhood inside the one from Proposition~\ref{prop:B.2}. Make it
small enough that the positive function \NSi{NS-F-P153-057}{\Phi} in the analytic axis profile has
no zeros. The formulas for \NSi{NS-F-P153-058}{Q_s,N_s}, computed from the reference profile, and
\eqref{eq:B.26} use only analytic parameter coefficients, their parameter derivatives, and forward
integrals, multiplied by cutoffs that depend only on \NSi{NS-F-P153-059}{y}. Their denominators do
not vanish on a smaller neighborhood. Exponentiation gives a nonzero
\NSi{NS-F-P153-060}{\phi}, so one neighborhood works across the whole fixed first collar. Radial
differentiation changes the bounds for the cutoffs, but it does not change the parameter domain.
\end{EESegment}

\begin{EESegment}{EE-TGT-P153-009}
Compactness and Proposition~\ref{prop:B.5} give an inner collar where the admissible stress cone
condition holds. Choose \NSi{NS-F-P153-061}{t_c} strictly inside this collar, and keep the resulting
rectangle free of all later changes. Neither \NSi{NS-F-P153-062}{t_c} nor the radial derivative
bounds are claimed to be uniform as \NSi{NS-F-P153-063}{\mathsf C} grows and the transition widths
shrink. The radial moment and pressure integrals run forward from the axis, so changes supported
later leave this rectangle unchanged.
\end{EESegment}
\end{proof}

\NSPage{154}%
\begin{EESegment}{EE-TGT-P154-001}
\subsection{Choosing a radial transition length before the amplitude.}\label{sec:B.7}
This subsection constructs the remaining transition \eqref{eq:B.34} to
\NSi{NS-F-P154-001}{E=P_*f(X/X_R)^{1/10}}, where
\NSi{NS-F-P154-002}{f=(1+\eta^2)^{-1}}. The radial moment functions are corrected in the next
subsection, in Proposition~\ref{prop:B.8}. First, Lemma~\ref{lem:B.7} and \eqref{eq:B.33} give a
bound for the logarithmic radial length \NSi{NS-F-P154-003}{T_{\mathrm{sh}}} that does not depend on
the amplitude \NSi{NS-F-P154-004}{\mathsf C} chosen later. Increasing that amplitude can then move
the matching radius \NSi{NS-F-P154-005}{X_R} outward without making the transition longer. This
makes the normalized discrepancies in the inner moments small, as shown in \eqref{eq:B.38} and
\eqref{eq:B.39}. The estimates below concern parameter derivatives. Bounds for derivatives in
\NSi{NS-F-P154-006}{X} do not have to be uniform in the radial transition widths. The number
\NSi{NS-F-P154-007}{\kappa_0\in(0,1/2)} in these bounds is the parameter that controls the shear
reduction in \eqref{eq:B.26}.
\end{EESegment}

\begin{lemma}\label{lem:B.7}
\begin{EESegment}{EE-TGT-P154-002}
Fix the outer-profile data, the axis data, \NSi{NS-F-P154-008}{\Lambda}, and a finite parameter
order \NSi{NS-F-P154-009}{k}. Set \NSi{NS-F-P154-010}{\ell_i=\log(\mathsf C E(X_i,\eta))} and
\NSi{NS-F-P154-011}{G_i=U(X_i,\eta)}. Let \NSi{NS-F-P154-012}{\omega_{\mathrm{fin}}} be the sum of
the logarithmic radial lengths of the final axial cutoff and the interpolation of
\NSi{NS-F-P154-013}{a} to \NSi{NS-F-P154-014}{.8} in the proof of
Proposition~\ref{prop:B.5}. The continuations built there satisfy
\end{EESegment}
\NSFormula{NS-F-P154-015}{display}{B.32}
\begin{equation}
\begin{aligned}
  \lVert\ell_i\rVert_{C_\eta^k}&\leq\mathcal B_k,\\
  \lVert G_i-U_*\rVert_{C_\eta^k}&\leq\frac{C_k^{\mathrm{nat}}}{\Lambda}
    +C_k^{\mathrm{join}}(\Lambda)(t_1+\kappa_0+\omega_{\mathrm{fin}}).
\end{aligned}
\tag{B.32}\label{eq:B.32}
\end{equation}
\begin{EESegment}{EE-TGT-P154-003}
The constant \NSi{NS-F-P154-016}{C_k^{\mathrm{nat}}} does not depend on either
\NSi{NS-F-P154-017}{\Lambda} or any sufficiently large \NSi{NS-F-P154-018}{\mathsf C}, once the
earlier axis data have been fixed. The bounds \NSi{NS-F-P154-019}{\mathcal B_k} and
\NSi{NS-F-P154-020}{C_k^{\mathrm{join}}(\Lambda)} may depend on
\NSi{NS-F-P154-021}{\Lambda}, but they do not depend on any sufficiently large
\NSi{NS-F-P154-022}{\mathsf C} or on any smaller radial transition widths.
\end{EESegment}
\end{lemma}
\begin{proof}
\begin{EESegment}{EE-TGT-P154-004}
Write \NSi{NS-F-P154-023}{L_0=\log(X_i/X_0)}. Let \NSi{NS-F-P154-024}{M_k,N_k} be the
amplitude-independent bounds for \NSi{NS-F-P154-025}{p_{1,\mathrm r},n_{s,\mathrm r}} from
Lemma~\ref{lem:B.4}. Before the last interpolation of \NSi{NS-F-P154-026}{a} to
\NSi{NS-F-P154-027}{.8}, the actual logarithmic slope is
\NSi{NS-F-P154-028}{-\kappa p_{1,\mathrm r}/2}, where
\NSi{NS-F-P154-029}{0<\kappa\leq1}. During that interpolation, the slope changes by a convex
interpolation to \NSi{NS-F-P154-030}{-.4}, and it stays constant after that. It follows that
\end{EESegment}
\NSFormula{NS-F-P154-031}{display}{none}
\[
  \lVert\log\phi(X_i)\rVert_{C_\eta^k}
  \leq\lVert\log\phi_*+\log\Phi(4,\mathord\cdot)\rVert_{C_\eta^k}
  +\tfrac12(M_k+.8)L_0.
\]
\begin{EESegment}{EE-TGT-P154-005}
Add \NSi{NS-F-P154-032}{\tfrac12\log(2X_i)} and use \eqref{eq:B.13}. This proves the first
claim. The first axial transition changes \NSi{NS-F-P154-033}{U} by at most
\NSi{NS-F-P154-034}{X_0e^{\bar t}N_kt_1/2}. The rest changes it by at most
\NSi{NS-F-P154-035}{X_iN_k\kappa_0/2}, with the same bound for the transition that takes this term
to zero. The approximation error for the analytic axis profile is
\NSi{NS-F-P154-036}{C_k^{\mathrm{nat}}/\Lambda}. Equation \eqref{eq:B.13} shows that its constant
does not depend on \NSi{NS-F-P154-037}{\Lambda}. These estimates prove \eqref{eq:B.32}. Only
integrals of bounded cutoff values were used, so no inverse power of a cutoff width appears.
\end{EESegment}
\end{proof}

\begin{EESegment}{EE-TGT-P154-006}
These bounds make it possible to choose the transition length without using the final amplitude.
Choose
\end{EESegment}
\NSFormula{NS-F-P154-038}{display}{B.33}
\begin{equation}
  T_{\mathrm{sh}}\geq20\lVert\sigma'\rVert_\infty
  \bigl(\mathcal B_0+\lVert\log f\rVert_\infty\bigr).
\tag{B.33}\label{eq:B.33}
\end{equation}
\begin{EESegment}{EE-TGT-P154-007}
For \NSi{NS-F-P154-039}{y_i=\log(X/X_i)}, prescribe
\end{EESegment}
\NSFormula{NS-F-P154-040}{display}{B.34}
\begin{equation}
  \log E=-\log\mathsf C+\frac{y_i}{10}
  +(1-\sigma(y_i/T_{\mathrm{sh}}))\ell_i
  +\sigma(y_i/T_{\mathrm{sh}})\log f,
  \qquad U=G_i.
\tag{B.34}\label{eq:B.34}
\end{equation}
\begin{EESegment}{EE-TGT-P154-008}
Then \NSi{NS-F-P154-041}{l=.6+T_{\mathrm{sh}}^{-1}\sigma'(y_i/T_{\mathrm{sh}})(\log f-\ell_i)}
lies in \NSi{NS-F-P154-042}{[.55,.65]}, so \NSi{NS-F-P154-043}{a\in[.7,.9]} and
\NSi{NS-F-P154-044}{b_s=0}.
\end{EESegment}

\begin{EESegment}{EE-TGT-P154-009}
Throughout the transition \eqref{eq:B.34}, the source still satisfies
\NSi{NS-F-P154-045}{S_q>1}. To check its sign, make all the changes before the final
\NSi{NS-F-P154-046}{.8} constant-slope interval small in the required
\NSi{NS-F-P154-047}{C_\eta^k} norms. The slope on that constant-slope interval does not depend on
the parameter. It follows that
\NSi{NS-F-P154-048}{\ell_i'=\widetilde\xi_0+(\log\Phi(4,\mathord\cdot))_\eta+}%
\NSPage{155}%
\NSi{NS-F-P155-001}{o(1)} and
\NSi{NS-F-P155-002}{G_i-U_*=O(\Lambda^{-1})+o(1)}. The old gradient satisfies
\end{EESegment}
\NSFormula{NS-F-P155-003}{display}{none}
\[
\begin{aligned}
  -H_c\ell_i'&\geq L\Lambda\chi-C\chi-C_{\sigma_*}\sqrt\chi-C/\Lambda-o(1),\\
  -H_c(\log f)'&=\frac{2\eta H_c}{1+\eta^2}\geq-Cj_0^2-C/\Lambda-o(1).
\end{aligned}
\]
\begin{EESegment}{EE-TGT-P155-001}
For the second line, use \NSi{NS-F-P155-004}{\eta H_*=(D+4d)\eta^2+dj_0\eta} and complete the
square. For the first line, use \eqref{eq:B.20}, which gives
\NSi{NS-F-P155-005}{H_*\partial_\eta\log\Phi=O(\chi)+O(\Lambda^{-1})}. Multiplying the
\NSi{NS-F-P155-006}{O(\Lambda^{-1})} error in the field by
\NSi{NS-F-P155-007}{\widetilde\xi_0} creates an error bounded by
\NSi{NS-F-P155-008}{C_{\sigma_*}\sqrt\chi}, and this error is absorbed in the same way as before.
Equation \eqref{eq:B.34} takes a convex combination of the two gradients. Since
\NSi{NS-F-P155-009}{-W} stays close to \NSi{NS-F-P155-010}{-W_*>2.8}, while
\NSi{NS-F-P155-011}{l\geq.55} and \NSi{NS-F-P155-012}{h,j_0} are small, it follows that
\NSi{NS-F-P155-013}{S_q>1}.
\end{EESegment}

\begin{EESegment}{EE-TGT-P155-002}
The inequality \NSi{NS-F-P155-014}{\mathcal D_Xp_1=XS_q/L-lp_1>0} holds at the barrier
\NSi{NS-F-P155-015}{p_1=2}, so a downward crossing is impossible. This preserves
\NSi{NS-F-P155-016}{p_1>2}. The whole transition, and the later interval of constant slope before
restoration, therefore continue to satisfy the strict inequalities that define the relaxed cone
condition.
\end{EESegment}

\begin{EESegment}{EE-TGT-P155-003}
\subsection{Matching the five radial moment functions exactly.}\label{sec:B.8}
The azimuthal profile \NSi{NS-F-P155-017}{E} now has the required reference form
\eqref{eq:A.7}. Equation \eqref{eq:B.32}, together with \eqref{eq:B.1}, says that the axial profile
\NSi{NS-F-P155-018}{U=G_i} is close to its reference value \NSi{NS-F-P155-019}{4\eta}. This
subsection matches the five cumulative radial integrals
\NSi{NS-F-P155-020}{M,I,J,S,C_p} from \eqref{eq:4.15}, so that the outer profile can be attached.
First choose the radial scale \NSi{NS-F-P155-021}{X_R} to make their normalized differences small.
Then cancel those differences with five fixed bumps while keeping the relaxed cone inequalities,
as stated in Proposition~\ref{prop:B.8}. Once the common pressure datum
\NSi{NS-F-P155-022}{\Pi_0} and all five moments match exactly, Lemma~\ref{lem:4.4} shows that every
later outer field stays the same.
\end{EESegment}

\begin{EESegment}{EE-TGT-P155-004}
Set
\end{EESegment}
\NSFormula{NS-F-P155-023}{display}{none}
\[
  X_R=X_i(\mathsf C P_*)^{10},
  \qquad
  x=X/X_R.
\]
\begin{EESegment}{EE-TGT-P155-005}
Also set \NSi{NS-F-P155-024}{X_{\mathrm{sep}}=X_ie^{T_{\mathrm{sh}}}}, the endpoint of the
transition \eqref{eq:B.34}. This endpoint is fixed before the amplitude is chosen.
\end{EESegment}

\begin{proposition}\label{prop:B.8}
\begin{EESegment}{EE-TGT-P155-006}
Take the amplitude \NSi{NS-F-P155-025}{\mathsf C} large enough, and then take the activation
transitions small enough. The profile built in Proposition~\ref{prop:B.5} and continued by
\eqref{eq:B.34} can then be continued to \NSi{NS-F-P155-026}{\log x=-5} while keeping the strict
relaxed cone condition. At that point, its fields, its pressure \NSi{NS-F-P155-027}{\Pi}, and all
five radial moment functions \NSi{NS-F-P155-028}{M,I,J,S,C_p} agree exactly with those of the
reference inner profile \eqref{eq:A.7}. The profile can then follow the reference outer radial
profile without changing either its prescribed pressure datum or any later value of
\NSi{NS-F-P155-029}{Q_s,N_s}.
\end{EESegment}
\end{proposition}
\begin{proof}
\begin{EESegment}{EE-TGT-P155-007}
First normalize the radial moment map so that its inverse bounds do not depend on the large radius
chosen later. Then show that the difference built up by the inner endpoint is small in these
normalized units, and cancel it with five fixed bumps.
\end{EESegment}

\begin{EESegment}{EE-TGT-P155-008}
Introduce \NSi{NS-F-P155-030}{\widehat H=\sqrt{2x}E}, and define the normalized radial moment
functions by
\end{EESegment}
\NSFormula{NS-F-P155-031}{display}{none}
\[
\begin{gathered}
  M=X_R\widehat M,\qquad
  I=X_R^{3/2}\widehat I,\qquad
  J=X_R^{3/2}\widehat J,\qquad
  S=X_R\widehat S,\\
  C_p=\displaystyle\int_0^x\frac{E(\xi,\eta)^2}{2\xi}\,d\xi,
  \qquad
  \Pi=\Pi_0+C_p.
\end{gathered}
\]
\begin{EESegment}{EE-TGT-P155-009}
Here \NSi{NS-F-P155-032}{\widehat M,\widehat I,\widehat J,\widehat S} are the integrals of
\NSi{NS-F-P155-033}{U,\widehat H,U\widehat H,U^2-E^2/2}, in that order, with respect to
\NSi{NS-F-P155-034}{x}. Substitute these definitions into \eqref{eq:4.16}. The result is
\end{EESegment}
\NSFormula{NS-F-P155-035}{display}{B.35}
\begin{equation}
\begin{aligned}
  W&=1-2D\eta\widehat M/x-d\widehat M_\eta/x,\\
  Q_s&=-W+\frac{(1-h)\widehat I-D\eta\widehat I_\eta-d\widehat J_\eta
                   +2(h-D)\eta\widehat J}{x\widehat H},\\
  N_s&=-WU+\frac{D(\widehat M-\eta\widehat M_\eta)+4h\eta\widehat S-d\widehat S_\eta}{x}
                   +4A\eta\Pi-d\Pi_\eta.
\end{aligned}
\tag{B.35}\label{eq:B.35}
\end{equation}
\NSPage{156}%
\begin{EESegment}{EE-TGT-P156-001}
In these formulas, \NSi{NS-F-P156-001}{X_R} no longer appears in
\NSi{NS-F-P156-002}{Q_s,N_s}. On the fixed interval
\NSi{NS-F-P156-003}{\mathcal R=[e^{-8},e^{-5}]}, the ideal functions
\NSi{NS-F-P156-004}{E_0=P_*fx^{1/10}} and \NSi{NS-F-P156-005}{Q_{s,0}} have positive minimum
values, and the ideal \NSi{NS-F-P156-006}{N_{s,0}} is bounded. These facts follow from the explicit
formula \eqref{eq:A.25}. On this interval, the map \eqref{eq:B.35} is Lipschitz in every chosen
finite parameter norm. Its inputs need one more parameter derivative, and its constants depend only
on the outer-profile data.
\end{EESegment}

\begin{EESegment}{EE-TGT-P156-002}
The same statements hold for the fixed restoration and moment-correction operations. To state the
allowed error in the cone inequalities, set
\NSi{NS-F-P156-007}{Q_{\min}=\min_{\mathcal R\times[-1,1]}Q_{s,0}>0},
\NSi{NS-F-P156-008}{e_*=\min E_0>0}, and
\NSi{NS-F-P156-009}{w_*=1+\max|N_{s,0}/(E_0Q_{s,0})|}. A tolerance depending only on the data of
the outer profile can make the following bounds hold throughout both operations.
\end{EESegment}
\NSFormula{NS-F-P156-010}{display}{B.36}
\begin{equation}
  Q_s\geq Q_{\min}/2,\qquad
  E\geq e_*/2,\qquad
  |a-.8|\leq.1,\qquad
  \left|\frac{N_s}{EQ_s}\right|\leq w_*,\qquad
  |b_s|\leq\frac{.1}{1+w_*}.
\tag{B.36}\label{eq:B.36}
\end{equation}
\begin{EESegment}{EE-TGT-P156-003}
These bounds give \NSi{NS-F-P156-011}{v_s\leq.9+.01/.7<1}. Set
\NSi{NS-F-P156-012}{G=Q_s-b_sN_s/(aE)}. Then
\end{EESegment}
\NSFormula{NS-F-P156-013}{display}{B.37}
\begin{equation}
  G\geq\frac67Q_s\geq\frac{3Q_{\min}}7,
  \qquad
  P_{\mathrm c}=\frac{X_Rx}{L}G.
\tag{B.37}\label{eq:B.37}
\end{equation}
\begin{EESegment}{EE-TGT-P156-004}
The inequality \NSi{NS-F-P156-014}{P_{\mathrm c}>2} follows by increasing
\NSi{NS-F-P156-015}{X_R} after the tolerance has been fixed. The tolerance for the moment correction does not have to
shrink as \NSi{NS-F-P156-016}{X_R} grows.
\end{EESegment}

\begin{EESegment}{EE-TGT-P156-005}
It remains to show that this fixed tolerance can be reached. In the normalized coordinate, the
endpoint \NSi{NS-F-P156-017}{X_{\mathrm{sep}}} is
\end{EESegment}
\NSFormula{NS-F-P156-018}{display}{B.38}
\begin{equation}
  x_{\mathrm{sep}}=X_{\mathrm{sep}}/X_R=e^{T_{\mathrm{sh}}}/(\mathsf C P_*)^{10}.
\tag{B.38}\label{eq:B.38}
\end{equation}
\begin{EESegment}{EE-TGT-P156-006}
On \NSi{NS-F-P156-019}{[0,X_{\mathrm{sep}}]}, the bounds already proved for
\NSi{NS-F-P156-020}{\log(\mathsf C E(X_i,\mathord\cdot))} and
\NSi{NS-F-P156-021}{U(X_i,\mathord\cdot)}, together with \eqref{eq:B.34}, bound
\NSi{NS-F-P156-022}{U} and \NSi{NS-F-P156-023}{\mathsf C E} in every required parameter norm.
Near zero, \NSi{NS-F-P156-024}{\mathsf C E=O(\sqrt X)}. The parts of the cumulative integrals that
come from \NSi{NS-F-P156-025}{0\leq x\leq x_{\mathrm{sep}}} therefore satisfy
\end{EESegment}
\NSFormula{NS-F-P156-026}{display}{B.39}
\begin{equation}
\begin{aligned}
  \lVert\widehat M\rVert_{C_\eta^k}+\lVert\widehat S\rVert_{C_\eta^k}
    &\leq C_kx_{\mathrm{sep}},\\
  \lVert\widehat I\rVert_{C_\eta^k}+\lVert\widehat J\rVert_{C_\eta^k}
    &\leq C_k\mathsf C^{-1}x_{\mathrm{sep}}^{3/2},\\
  \lVert C_p\rVert_{C_\eta^k}
    &\leq C_k\mathsf C^{-2}\qquad(x=x_{\mathrm{sep}}).
\end{aligned}
\tag{B.39}\label{eq:B.39}
\end{equation}
\begin{EESegment}{EE-TGT-P156-007}
The matching integrals for the ideal profile have positive weights, and they also go to zero as
positive powers of \NSi{NS-F-P156-027}{x_{\mathrm{sep}}}. At
\NSi{NS-F-P156-028}{x_{\mathrm{sep}}}, the pressure increment of the reference profile is
\NSi{NS-F-P156-029}{(5/2)P_*^2f^2x_{\mathrm{sep}}^{1/5}=O(\mathsf C^{-2})}. The constants in all
these vanishing errors may depend on the axis and shape parameters, which have already been fixed.
\end{EESegment}

\begin{EESegment}{EE-TGT-P156-008}
After \NSi{NS-F-P156-030}{X_{\mathrm{sep}}}, equation \eqref{eq:B.34} gives exactly the ideal
angular profile because
\NSi{NS-F-P156-031}{\mathsf C^{-1}f(X/X_i)^{1/10}=P_*fx^{1/10}}. Choose the amplitude large
enough that \NSi{NS-F-P156-032}{x_{\mathrm{sep}}<e^{-8}}. On
\NSi{NS-F-P156-033}{-8<\log x<-7}, bring \NSi{NS-F-P156-034}{G_i} smoothly back to
\NSi{NS-F-P156-035}{4\eta}, while keeping \NSi{NS-F-P156-036}{E} equal to its ideal value. The
difference \NSi{NS-F-P156-037}{G_i-4\eta} can be made as small as needed by first choosing
\NSi{NS-F-P156-038}{j_0}, then choosing \NSi{NS-F-P156-039}{\Lambda}, and finally choosing the
small transitions in \eqref{eq:B.32}.
\end{EESegment}

\begin{EESegment}{EE-TGT-P156-009}
Integrating over the rest of the fixed interval gives bounds depending only on the outer-profile data,
because all the weights involved are integrable positive powers at zero. At the point where the
moments are corrected, the five normalized moment differences are therefore at most
\NSi{NS-F-P156-040}{C_k\lVert G_i-4\eta\rVert_{C_\eta^{k+1}}+o_{\mathrm C}(1)}. The same estimate
controls the field values, their logarithmic radial derivatives, and
\NSi{NS-F-P156-041}{Q_s,N_s} on the fixed restoration patch.
\end{EESegment}

\begin{EESegment}{EE-TGT-P156-010}
On \NSi{NS-F-P156-042}{-6<\log x<-5}, add two fixed radial bumps to
\NSi{NS-F-P156-043}{U} and three fixed radial bumps to \NSi{NS-F-P156-044}{E}. At the ideal
profile, replace the \NSi{NS-F-P156-045}{J} row by \NSi{NS-F-P156-046}{J-4\eta I}, and replace
the \NSi{NS-F-P156-047}{S} row by \NSi{NS-F-P156-048}{S-8\eta M}. Apart from fixed nonzero
normalizing factors, the derivative of the normalized moment map has the following two blocks.
\end{EESegment}
\NSFormula{NS-F-P156-049}{display}{none}
\[
  U\ \text{weights}:\quad 1,\ fx^{3/5};
  \qquad
  E\ \text{weights}:\quad x^{1/2},\ fx^{1/10},\ fx^{-9/10}.
\]
